\documentclass[a4paper]{article}

\usepackage{anyfontsize}
\usepackage{amsmath,amssymb,mathrsfs}
\usepackage{autobreak}
\allowdisplaybreaks
\usepackage{bm}
\usepackage{indentfirst}
\setlength{\parindent}{0em}
\usepackage{parskip}
\usepackage{geometry}
\geometry{top=3cm,bottom=3cm,left=2.75cm,right=2.75cm}

\begin{document}

\title{\textbf{Exercise 1}}
\author{}
\date{}

\maketitle

\subsection*{Problem 1: Measuring flat universes}

In the following consider two flat $(k=0)$, expanding universes: one that is matter dominated, i.e. $\Omega_\mathrm{m}=1$, and one that is dominated by a cosmological constant, i.e. $\Omega_\Lambda=1$.

\begin{itemize}
    \item[(a)] Derive for both universes an expression for the comoving distance
    \[
    D(z) = \omega(z)R_0
    \]
    as a function of redshift $z$.

    \item[(b)] The "observable universe" is the part of the universe up to redshift $z=\infty$. Compute for both of our universes the comoving volume of this region, i.e. the volume of the corresponding region on the spatial hypersurface of constant time $\tau_0$. Is there a qualitative difference between the two volumes? What is the meaning of the "observable universe" in the context of causality?
    
    \textit{Hint:} The spatial volume element for the Robertson Walker (RW) metric at time $\tau$ is
    \[
    \mathrm{d}V = R^3(\tau)S_k^2(\omega)\sin\theta\,\mathrm{d}\theta\,\mathrm{d}\phi\,\mathrm{d}\omega
    \]
\end{itemize}

\subsection*{Problem 2: Change of redshift with time}

The redshift of a given source of light should change with time and in principle we could detect this effect by observing the source over a certain timespan. Assume a flat concordance cosmology with $H_0=73\,\mathrm{km/(s\cdot Mpc)}$, $\Omega_\mathrm{m}=0.3$ and $\Omega_\Lambda=0.7$. How much would the redshift of a source at redshift $z=1$ change within ten years? Do you think this might be observable in practice and what complications might arise?

\end{document}